\begin{ExerciseBlock}
\ExerciseHeading{02}{Absolute und relative Messunsicherheit}{laboratory_measurements_absolute_relative_uncertainty}
Geben Sie die Messwerte mit der jeweils anderen Form der Unsicherheit an.
\begin{enumerate}
\item Die Masse beträgt $m=\qty{50.0}{\gram}$ mit $u(m)=\qty{0.2}{\gram}$. Geben Sie $m$ mit relativer Unsicherheit $u(m)/m$ an.
\item Die Länge beträgt $L=\qty{1.5}{\metre}$ mit relativer Unsicherheit von \qty{2}{\percent}. Geben Sie $L$ mit absoluter Unsicherheit $u(L)$ an.
\end{enumerate}
\end{ExerciseBlock}
\begin{Solution}
\par\medskip\noindent\textbf{Musterlösung}\par\nopagebreak[4]
\begin{enumerate}
\item Die relative Unsicherheit ist $u(m)/m=\qty{0.2}{\gram}/\qty{50.0}{\gram}=0{,}004=\qty{0.4}{\percent}$. Damit ist $m=\qty{50.0}{\gram}(1\pm0{,}004)$.
\item Die absolute Unsicherheit ist $u(L)=L\cdot0{,}02=\qty{0.03}{\metre}$. Damit ist $L=(1{,}50\pm0{,}03)\,\unit{\metre}$.
\end{enumerate}
\end{Solution}
